\documentclass{article}
\usepackage{hyperref}
\hypersetup{colorlinks=true, linkcolor=blue, urlcolor=magenta}

\title{Hyperlinks and Footnotes}
\author{platex demo}
\date{}

\begin{document}
\maketitle

\section{External links}
The platex library targets Overleaf-equivalent output; see
\url{https://overleaf.com} for comparison, or use a
\href{https://overleaf.com}{custom link label}.

\section{Internal cross-references}
This section is labeled for cross-referencing purposes\footnote{Footnotes
render at the bottom of the page and are automatically numbered.}, and can
be jumped to via Section~\ref{sec:target} below.

\section{Target section}
\label{sec:target}
This is the section referenced above, demonstrating that \texttt{hyperref}
turns \texttt{\textbackslash ref} output into a clickable internal link in
the resulting PDF.

Another footnote here\footnote{Footnotes can appear anywhere in the flow of
text, not just at section boundaries.} to show consecutive numbering.

\end{document}
